\chapter[Oscillators and stochastic dynamics]{Oscillators, bifurcations, and stochastic dynamics}
\label{ch:oscillator-foundations}
Oscillatory-looking trajectories can emerge from a dynamical
system, an externally driven filter, or chance structure in
a finite recording. These are mathematically distinct accounts.

\section{Hopf normal form and an actual periodic attractor}
The complex-valued ordinary differential equation
$\dot z=(\mu+i\omega)z-|z|^2z$ with $\omega\ne0$
can be written $z=re^{i\theta}$, giving
\[
\dot r=\mu r-r^3,\qquad \dot\theta=\omega.
\]
For $\mu<0$ the origin attracts nearby trajectories.
For $\mu>0$ the origin repels, while the circle
$r=\sqrt\mu$ is an attracting periodic orbit. Its period
is $2\pi/|\omega|$.
This explicit derivation demonstrates a supercritical Hopf-type
normal form; deciding whether experimental LFP oscillations
arise from an actual Hopf bifurcation requires more evidence.

\section{Phase reduction and weak coupling}
A stable limit-cycle oscillator with phase coordinate $\theta$
responds to weak perturbation $u(t)$ to leading order as
\[
\dot\theta=\omega+Z(\theta)\cdot u(t),
\]
where $Z$ is the infinitesimal phase-response curve (PRC).
For two weakly coupled oscillators, averaging can produce
a phase-difference equation
$\dot\psi=\Delta\omega+\varepsilon H(\psi)$.
Locking requires a stable fixed point of this equation;
oscillatory phase locking need not imply anatomical
synapses, and non-sinusoidal harmonics can generate apparent PAC.

\section{jPCA and skew-symmetric dynamics}
If $\dot x=Mx$ for a real skew-symmetric matrix
$M^\top=-M$, then
\[
\frac{d}{dt}\|x\|_2^2=2x^\top Mx=0.
\]
A real skew-symmetric matrix has orthogonal real normal form
consisting of planar blocks
$\begin{psmallmatrix}0&-\omega_j\\\omega_j&0\end{psmallmatrix}$
(and possibly zeros); its nonzero eigenvalues occur as
imaginary pairs $\pm i\omega_j$. Thus a fitted rotational
component conserves Euclidean energy in the ideal autonomous
model. The file \path{jpca_dynamics.m}
fits such a skew generator to finite-difference derivatives
\cite{churchland2012},
not a complete state-space model or a proof of physiological cycles.

\section{Ornstein--Uhlenbeck noise and the Fokker--Planck equation}
For an It\^o stochastic differential equation
$dX_t=a(X_t)dt+b(X_t)dW_t$, a twice differentiable
function $f$ obeys It\^o's formula
\[
df(X_t)=\left(af'+\tfrac12 b^2f''\right)dt
+b f'\,dW_t.
\]
Its density, under regularity and boundary assumptions,
evolves by the Fokker--Planck equation
$\partial_tp=-\partial_x(ap)
 +\tfrac12\partial_{xx}(b^2p)$.
For the Ornstein--Uhlenbeck process
$dX_t=-\kappa X_tdt+\sigma dW_t$, with $\kappa>0$,
the stationary variance is $\sigma^2/(2\kappa)$ and the
stationary autocovariance is
$\sigma^2e^{-\kappa|\tau|}/(2\kappa)$.
An exponential autocorrelation is therefore compatible with
a stochastic relaxation process and need not imply an intrinsic
neuronal pacemaker or periodic oscillator.

\begin{example}[Simulation check]
Simulate a damped OU fluctuation and a noisy Hopf limit cycle
at the same sampling frequency. Compare autocorrelation,
PSD and phase stability under different recording lengths.
A narrow spectral peak alone does not establish the Hopf model.
\end{example}
% --- expanded edition (October 2026) ---
\section{Locking to a periodic input: the Adler equation}
Let an oscillator with natural frequency $\omega$ be driven weakly at
frequency $\Omega$ with a sinusoidal phase interaction. Averaging gives
for the phase difference $\psi=\theta-\Omega t$
\[
\dot\psi=\Delta\omega-\varepsilon\sin\psi,\qquad\Delta\omega=\omega-\Omega.
\]
Fixed points exist iff $|\Delta\omega|\le\varepsilon$, at
$\sin\psi^\ast=\Delta\omega/\varepsilon$; the one with $\cos\psi^\ast>0$ is
stable. Inside this range the oscillator is \emph{entrained}: it runs
at the driving frequency with a constant phase lag that depends on the
detuning. Outside it, $\psi$ drifts, with mean drift frequency
$\sqrt{\Delta\omega^2-\varepsilon^2}$ (phase slips separated by long
quasi-locked epochs). Plotted in the $(\Omega,\varepsilon)$ plane, the
locking region is a wedge, the first \emph{Arnold tongue}; higher-order
tongues for $p\!:\!q$ locking appear with nonsinusoidal interaction
functions.

The equation predicts what to look for in rhythmic stimulation
experiments: entrainment should hold over a range of stimulation
frequencies that grows with intensity, the phase lag should vary
systematically across that range, and the effect should vanish abruptly
outside it. A response that simply follows every stimulus frequency with
constant lag is more consistent with a driven filter than with an
entrained oscillator.

\section{The Kuramoto model}
For $N$ oscillators with natural frequencies $\omega_j$ drawn from a
density $g$ and all-to-all coupling,
\[
\dot\theta_j=\omega_j+\frac KN\sum_{k=1}^N\sin(\theta_k-\theta_j).
\]
With the order parameter $re^{i\Phi}=\frac1N\sum_ke^{i\theta_k}$ the
system becomes $\dot\theta_j=\omega_j+Kr\sin(\Phi-\theta_j)$: each
oscillator is driven by the mean field, an Adler equation with
$\varepsilon=Kr$. Oscillators with $|\omega_j-\Omega|\le Kr$ lock;
the others drift. Self-consistency of $r$ gives Kuramoto's result: for a
symmetric unimodal $g$, the incoherent state $r=0$ loses stability at
\[
K_c=\frac{2}{\pi g(0)},
\]
and for a Lorentzian $g$ of half-width $\gamma$, $K_c=2\gamma$ and
$r=\sqrt{1-K_c/K}$ for $K>K_c$. The order parameter $r$ is the
population analogue of the phase-locking value, and the model shows that
partial synchrony, a fraction of locked units with others drifting, is a
generic state rather than a sign of measurement failure.

\section{Center manifolds and the Hopf normal form}
The normal form of this chapter is not an arbitrary model: near a
bifurcation, it is what any smooth system reduces to.
\begin{theorem}[Center manifold, stated]
Let $\dot x=f(x,\mu)$ with $f(0,\mu)=0$, and suppose $Df(0,0)$ has a pair
of eigenvalues $\pm i\omega_0$ and all others with negative real part.
Then for $\mu$ near $0$ there is a locally invariant two-dimensional
manifold through the origin that attracts nearby trajectories, and the
dynamics on it are smoothly equivalent to
$\dot z=(\alpha(\mu)+i\omega(\mu))z+\ell_1|z|^2z+O(|z|^4)$.
\end{theorem}
The sign of the real part of the \emph{first Lyapunov coefficient}
$\ell_1$ decides the type: negative gives the supercritical Hopf of
this chapter (small limit cycles growing as $\sqrt\mu$); positive gives
a subcritical Hopf, in which a large-amplitude oscillation appears
abruptly and coexists with the resting state (bistability, hysteresis).
Neuron models distinguish these: class-2 excitability with a
subcritical Hopf begins firing at a nonzero frequency and can show
bistability between rest and repetitive firing.

\section{Noise-driven oscillations near a Hopf point}
Below the bifurcation ($\mu<0$), noise $dz=(\mu+i\omega)z\,dt+\sigma\,dW$
(complex $W$) produces a stationary Gaussian process whose spectrum is a
Lorentzian peaked at $\omega/(2\pi)$ with half-width $|\mu|/(2\pi)$ in
hertz. Its amplitude has a Rayleigh distribution and its phase diffuses.
Above the bifurcation, the amplitude concentrates near $\sqrt\mu$ and the
phase still diffuses, with a much narrower peak. The two regimes can
produce similar power spectra; the amplitude distribution and the
duration of phase coherence discriminate them better than the spectrum.

\section*{Exercises for this chapter}
\begin{exercise}
Solve the Adler equation outside the locking range and derive the mean
drift frequency $\sqrt{\Delta\omega^2-\varepsilon^2}$.
\end{exercise}
\begin{exercise}
For a Lorentzian $g$ with $\gamma=1$\,rad/s, compute $K_c$ and the order
parameter at $K=4$.
\end{exercise}
\begin{exercise}
Show that the OU process $dX=-\kappa X\,dt+\sigma\,dW$ has spectral
density proportional to $1/(\kappa^2+(2\pi f)^2)$.
\end{exercise}
\begin{exercise}
For $\dot r=\mu r+r^3-r^5$, find the equilibria and their stability, and
describe the hysteresis as $\mu$ crosses zero.
\end{exercise}
